% !TEX root = ../../tree_notes.tex
% =====================================================================
%  9.3  Range Minimum Query via the Cartesian Tree   sec:c26-rmq
%  Exports: thm:c26-rmqlca  lem:c26-cartmin  fig:c26-rmqfig
%           lst:c26-rmq  ex:c26-rmq-run
%           prob:c26-rmq1647  lst:c26-rmq1647-soln
% =====================================================================
\section{Range Minimum Query via the Cartesian Tree}\label{sec:c26-rmq}

\begin{previouslybox}
\begin{enumerate}[nosep, leftmargin=1.6em]
  \item What are the two orders of a Cartesian tree, and which is the root? (\defref{c26-cartesian})
  \item \lemref{c26-span}: node $i$'s subtree is the window on which $a[i]$ is the minimum. If a range
        $[l,r]$ has its minimum at position $m$, where do $l$ and $r$ sit relative to $m$'s window?
  \item In a heap, comparing a node to its ancestor: is the ancestor's value smaller or larger? (smaller --
        min-heap)
  \item Which structure from Chapter~16 (\chref{c11}) answers "lowest common ancestor of two nodes" after
        $\bigO{n\log n}$ preprocessing? (binary lifting)
\end{enumerate}
\end{previouslybox}

\begin{topicintro}
Here the Cartesian tree stops being a metaphor and becomes an \emph{index}. You are handed a fixed array
and many questions of the form ``what is the minimum of $a[l..r]$?'' -- the \term{range minimum query},
the static sibling of the segment tree. The punchline of this section is a one-line identity:

\begin{center}
\setlength{\fboxsep}{6pt}\fbox{\normalsize$\displaystyle\min\{a[l],\dots,a[r]\} \;=\; a\big[\lca(l,r)\big]$
\quad in the min-Cartesian tree.}
\end{center}

The minimum of any contiguous range is the \emph{value stored at the lowest common ancestor} of the two
endpoints, when the endpoints are read as nodes of the tree. Proving it is a three-line consequence of
\lemref{c26-span}; using it turns RMQ into an LCA call, and (this is the deeper equivalence Chapter~16
plays out) LCA is itself an RMQ on the Euler tour -- the two problems are the same problem in two
languages. We state and prove the reduction, build the tree, and answer queries with binary lifting,
explicitly \emph{handing the LCA engine over} to Chapter~16 rather than re-deriving it.
\textbf{Prerequisites.} \secref{c26-cartesian} (the span lemma is the entire proof); a black-box LCA
(binary lifting, Chapter~16, \chref{c11}).

\textbf{After this section you can}
\begin{itemize}[nosep]
  \item prove that a range minimum equals the value at the LCA of its endpoints in the Cartesian tree; \hfill\lv{3}--\lv{4}
  \item answer $q$ static RMQs in $\bigO{(n+q)\log n}$ by combining the $\bigO{n}$ build with lifted LCA; \hfill\lv{4}
  \item explain the two-way RMQ $\equiv$ LCA equivalence and say which chapter owns which direction. \hfill\lv{4}
\end{itemize}
\end{topicintro}

\begin{recognitionbox}[Pattern card -- static range extremes]
\reckey{a \emph{fixed} array, \emph{no} updates, and many "min/max over $[l,r]$" questions -- the offline
regime where the Cartesian tree doubles as a query index.}
\recuse{
  \item "minimum in range $[a,b]$", many queries, array never changes -- CSES 1647 (\probref{c26-rmq1647});
  \item "LCA of two nodes" dressed as an array problem, or vice-versa -- the equivalence is the tool;
  \item "is $a[m]$ the min of this window?" -- check $m$ against $\lca$ of the endpoints' nodes.
}
\recnotuse{
  \item the array is \emph{updated} between queries -- you need a segment tree (Chapter~18--19,
        \chref{c14}/\chref{c13}); the Cartesian tree is frozen at build time;
  \item a \emph{single} whole-array min -- one scan;
  \item "sum over a range" -- the min-heap order is what makes LCA work; sums have no such reduction, use
        a prefix array or segment tree.
}
One-line test: \emph{static array, range min, lots of queries -- and you would like to think in terms of
a tree? Then $\min = $ the value at an LCA, and the reduction is the whole solution.}
\end{recognitionbox}

\subsection{Intuition}

Take $A=(3,1,4,2)$ again. Ask for the minimum of the range covering positions $2$ through $3$: $\min(1,4)=1$,
attained at position $2$. Now look at the Cartesian tree \figref{c26-cartA}: the node for position $2$
(value $1$) is the \emph{root}, and positions $2$ and $3$ are both in its subtree, so their lowest common
ancestor is exactly position $2$ -- the minimum. Every range tells the same story: the position of its
smallest value is the node from which both endpoints are reachable on the way up, and no lower node works
because a lower node would have a \emph{larger} value (the heap points up toward smaller).

Why can there be no common ancestor below the minimum? In a min-heap the values only get \emph{smaller}
going up, so a deeper common ancestor would have to be smaller than the range's own minimum -- impossible,
since it stores a value from inside the range. Why is the minimum \emph{an} ancestor of both endpoints?
Because \lemref{c26-span} says the minimum's subtree is its dominance window, and every value in $[l,r]$
is $\ge$ that minimum, so $[l,r]$ lies entirely inside the minimum's window -- meaning both $l$ and $r$
are in its subtree. Ancestor of both, none deeper possible: it \emph{is} the LCA.

\begin{figure}[H]
\centering
\begin{tikzpicture}[diagram, scale=0.92, transform shape]
  \node[troot] (n2) at (0,0)  {$a_2{=}1$};
  \node[tleaf] (n1) at (-2,-1.5) {$a_1{=}3$};
  \node[tnode] (n4) at (2,-1.5)  {$a_4{=}2$};
  \node[tleaf] (n3) at (1,-3.0)  {$a_3{=}4$};
  \draw[tedgehl] (n2) -- (n1); \draw[tedgehl] (n2) -- (n4); \draw[tedgehl] (n4) -- (n3);
  % query [1,3]: endpoints a1 and a3, LCA a2
  \node[draw=warningred, thick, circle, inner sep=1pt, minimum size=10mm] at (n1) {};
  \node[draw=warningred, thick, circle, inner sep=1pt, minimum size=10mm] at (n3) {};
  \node[annot, align=left, font=\scriptsize] at (0,1.2)
    {query $\min\{a_1..a_3\}$: endpoints $1,3$ (red)};
  \node[annot, align=left, font=\scriptsize] at (-5.0,-2.2)
    {\textbf{$\lca(1,3)=2$}\\ $a[\lca]=1$\\ $=\min(3,1,4)$};
  \draw[flow, dashed] (n1) -- (n2); \draw[flow, dashed] (n3) -- (n4) -- (n2);
\end{tikzpicture}
\caption{The reduction on $A=(3,1,4,2)$: to find $\min\{a_1,a_2,a_3\}$ read the two endpoint nodes $1$
and $3$ (ringed), climb to their lowest common ancestor -- node $2$ -- and take its value, $1$. That is
$\thmref{c26-rmqlca}$; the dashed paths are the way up to the shared ancestor.}\label{fig:c26-rmqfig}
\end{figure}

\subsection{Definitions and Notation}

\begin{lemma}{The range minimum is the shared ancestor}{c26-cartmin}
Let $m$ be a position attaining the minimum of $a[l..r]$ in the min-Cartesian tree. Then $m$ is an
ancestor of both node $l$ and node $r$, and no strict descendant of $m$ is a common ancestor of $l$ and
$r$. Hence $\lca(l,r)=m$.
\end{lemma}
\begin{proofbox}
By \lemref{c26-span}, $m$'s subtree spans its dominance window $(L_m,R_m)$. Every element of $[l,r]$ is
$\ge a[m]$ (as $m$ is the range minimum), so no element of $[l,r]$ is a strictly smaller value that could
sit at or beyond $L_m,R_m$; therefore $[l,r]\subseteq(L_m,R_m)$ and $l,r$ both lie in $m$'s subtree --
i.e.\ $m$ is a common ancestor of $l,r$. For the second claim, $m$'s left subtree holds exactly the
indices $<m$ and its right subtree the indices $>m$ (BST-on-index order); a descendant of $m$ lies in one
of these subtrees and so can contain indices only from one side of $m$, never both a left endpoint
$l\le m$ and a right endpoint $r\ge m$ -- unless $m$ equals one endpoint, in which case $m$ is trivially
the LCA. Thus no proper descendant is common to $l$ and $r$, and $m$ is the \emph{lowest} common ancestor. \end{proofbox}

\begin{theorem}{RMQ reduces to LCA on the Cartesian tree}{c26-rmqlca}
Build the min-Cartesian tree of a static array $a[1..n]$. For every $1 \le l \le r \le n$,
$\min\{a[l],\dots,a[r]\} = a\big[\lca(l,r)\big]$, where $\lca$ is taken over the node set
$\{1,\dots,n\}$ of that tree. Consequently a range-minimum query is answered in the time of a single
lowest-common-ancestor query, and any LCA engine -- binary lifting (Chapter~16, \chref{c11}),
Euler-tour plus sparse table, or Farach-Colton--Bender -- becomes an RMQ engine for $a$.
\end{theorem}
\begin{proofbox}
Directly from \lemref{c26-cartmin}: the position $m=\arg\min$ over $[l,r]$ satisfies
$\lca(l,r)=m$, so the queried minimum $a[m]=a[\lca(l,r)]$. The identity holds for the endpoints $l,r$
regardless of the interior, so a static $\lca$ oracle answers the whole family of queries; building the
tree (\coderef{c26-stackbuild}) is $\bigO{n}$, so the total cost is (tree build) $+$ (LCA preprocessing)
$+$ (per-query LCA), exactly as the chosen LCA engine prices them. \end{proofbox}

\begin{notebox}[The other direction: LCA $\equiv$ RMQ]
The reduction runs both ways, and that is the point of the equivalence Chapter~16 makes precise: given a
tree, list its nodes by \emph{Euler-tour entry time}; the LCA of $u,v$ is the node of \emph{minimum depth}
between $\tin(u)$ and $\tin(v)$ in that tour -- an RMQ on the depth array. So "static array min over a
range" and "lowest common ancestor in a tree" are mutually reducible. This section uses
array$\to$LCA (build a tree, query it); Chapter~16 (\chref{c11}) uses tree$\to$array (Euler tour, then
RMQ) and owns the LCA machinery -- binary lifting, and the Euler-tour/sparse-table route -- that we call
as a black box. If you only ever need \emph{array} RMQ, the sparse table (Chapter~19, \chref{c13}) answers
each query in $\bigO{1}$ with $\bigO{n\log n}$ build and never builds a tree; the Cartesian reduction
matters for the conceptual bridge and when you already \emph{have} the tree.
\end{notebox}

\subsection{The algorithm: build the tree, then lift}

The pipeline is exactly \coderef{c26-stackbuild} feeding an LCA oracle. Record each node's parent during
the build, run one depth pass, fill the binary-lifting table, and answer each query by lifting
$\lca(l,r)$ and printing $a$ of that node.

\begin{lstlisting}[caption={Array RMQ via the Cartesian tree: stack build records parents, binary lifting answers LCA, the query is a[LCA]}, label={lst:c26-rmq}]
// a[1..n]; q queries (l, r); outputs min of a[l..r] via a[LCA(l,r)] in the min-Cartesian tree.
vector<int> par(n+1,0), Lc(n+1,0), Rc(n+1,0), st;
for (int i = 1; i <= n; ++i) {                     // the O(n) Cartesian build (lst:c26-stackbuild)
    int last = 0;
    while (!st.empty() && a[st.back()] > a[i]) { last = st.back(); st.pop_back(); }
    if (!st.empty()) { Rc[st.back()] = i; par[i] = st.back(); }
    if (last)      { Lc[i] = last; par[last] = i; }
    st.push_back(i);
}
int root = st.front();
vector<int> dep(n+1, 0);                            // depth, iterative from the root (no recursion)
vector<int> fs = {root};
while (!fs.empty()) {
    int u = fs.back(); fs.pop_back();
    if (Lc[u]) { dep[Lc[u]] = dep[u] + 1; fs.push_back(Lc[u]); }
    if (Rc[u]) { dep[Rc[u]] = dep[u] + 1; fs.push_back(Rc[u]); }
}
int LOG = 1; while ((1 << LOG) < n) ++LOG;
vector<vector<int>> up(LOG, vector<int>(n+1, 0));
for (int v = 1; v <= n; ++v) up[0][v] = par[v];
for (int k = 1; k < LOG; ++k)
    for (int v = 1; v <= n; ++v) up[k][v] = up[k-1][ up[k-1][v] ];
auto lca = [&](int u, int v) {                      // standard binary-lifting LCA
    if (dep[u] < dep[v]) swap(u, v);
    int d = dep[u] - dep[v];
    for (int k = 0; k < LOG; ++k) if (d >> k & 1) u = up[k][u];
    if (u == v) return u;
    for (int k = LOG - 1; k >= 0; --k)
        if (up[k][u] != up[k][v]) { u = up[k][u]; v = up[k][v]; }
    return par[u];
};
for (int qi = 0; qi < q; ++qi) { int l, r; /* read */ cout << a[ lca(l, r) ] << '\n'; }
\end{lstlisting}
\begin{complexitybox}
Time $\bigO{n\log n + q\log n}$: $\bigO{n}$ to build, $\bigO{n\log n}$ for the lifting table, $\bigO{\log
n}$ per query. Memory $\bigO{n\log n}$ for \code{up}. Verdict: correct and comfortably inside $n,q\le
2\cdot10^5$. If you will \emph{only ever} answer array RMQ, the sparse table gives $\bigO{1}$ per query at
the same $\bigO{n\log n}$ build (Chapter~19, \chref{c13}); the Cartesian route is the one to know when the
tree is already part of your structure or you want the LCA bridge.
\end{complexitybox}

\subsection{Worked Example}

\begin{example}{Three queries on $A=(3,1,4,2)$}{c26-rmq-run}
Tree \figref{c26-cartA}: root $a_2{=}1$, left $a_1$, right $a_4$ with left child $a_3$.\\
\emph{$[1,3]$:} $\lca(1,3)$. Node $1$'s ancestors: $\{1,2\}$; node $3$'s: $\{3,4,2\}$; lowest common $=2$,
value $1$. Indeed $\min(3,1,4)=1$.\\
\emph{$[3,4]$:} $\lca(3,4)$. Node $4$ is the parent of node $3$, so $\lca=4$, value $2$;
$\min(4,2)=2$.\\
\emph{$[1,4]$:} $\lca(1,4)=2$ (the root), value $1=\min(3,1,4,2)$. Each query is one LCA, and each LCA is
the range's minimum -- \thmref{c26-rmqlca}. \checkmark
\end{example}

\subsection{Common Pitfalls}\label{sec:c26-rmq-pitfalls}

\begin{warningbox}[Trap 1 -- taking the LCA of values instead of positions]
The tree's nodes \emph{are the indices} $1..n$; you query $\lca(l,r)$ on the \emph{endpoint positions}
and then read $a$ of the result. Passing the values $a[l],a[r]$ as node names is a category error: the
heap order is on values, but the tree structure (parent/child, and hence LCA) is on the index-labeled
nodes.
\end{warningbox}

\begin{warningbox}[Trap 2 -- a recursive depth pass on a skewed tree]
A sorted array builds a Cartesian tree of depth $n$; a recursive DFS to fill \code{dep} then overflows
the call stack (RE) at $n=2\cdot10^5$. Use the explicit-stack sweep in \coderef{c26-rmq}, or compute
depth during the build. The same caution that haunts any pointer tree on adversarial input, back in
Chapter~7 (\chref{c05}).
\end{warningbox}

\begin{warningbox}[Trap 3 -- reinventing LCA you were given]
Do not hand-roll the Euler-tour-plus-sparse-table or the lifting table from scratch in an RMQ problem if
Chapter~16 already supplies it -- that is the chapter's anti-repetition contract. Build the Cartesian
tree, look up $\lca$ through the provided engine, and output $a[\lca]$. The \emph{reduction} is this
section's contribution; the \emph{LCA} is not.
\end{warningbox}

\subsection{Problem Lab}

The anchor is CSES 1647 -- Static Range Minimum Queries -- the purest statement of the reduction: a
frozen array, $q$ ranges, print each minimum. We answer every query as $a[\lca]$.

% ---------------------------------------------------------------------------
\solvedproblem{c26-rmq1647}{Static Range Minimum Queries}{CSES 1647}
  {https://cses.fi/problemset/task/1647}{CSES · L3}
  {Cartesian tree build, binary-lifting LCA, RMQ equals value at LCA}

\begin{problemstatement}[Static Range Minimum Queries]
Given an array of $n$ integers, process $q$ queries of the form "what is the minimum value in range
$[a,b]$?". The array never changes. $1 \le n,q \le 2\cdot10^{5}$, $1 \le x_i \le 10^{9}$, and
$1 \le a \le b \le n$ (1-indexed). Print the answer to each query on its own line.
\end{problemstatement}

\begin{attemptgate}[Attempt it yourself first -- \lv{4}, give yourself 30 minutes]
You may already know the sparse-table answer; force yourself to answer it \emph{as the reduction} so the
bridge is real, not decorative.
\peek{Build the min-Cartesian tree of the array. For a query $[a,b]$, which single node holds the answer?
(\thmref{c26-rmqlca}: the LCA of $a$ and $b$.)}
\nudge{You need LCA, so record parents during the build, get depths, and answer LCA by binary lifting.
Then print $a[\lca(a,b)]$. Preprocess $\bigO{n\log n}$, per query $\bigO{\log n}$.}
\fullhint{Exactly \coderef{c26-rmq}: stack build for \code{par/Lc/Rc}, iterative depth, a \code{up[k][v]}
lifting table, standard lift-to-LCA, output $a[\lca]$. On the sample the four answers are $2,1,1,4$.}
\end{attemptgate}

\solstep{Understanding the Problem}
Static array, no updates, up to $2\cdot10^5$ range-min queries. A per-query linear scan is $\bigO{nq}
\approx 4\cdot10^{10}$ -- far too slow. The standard engine is a sparse table; but this section's lens is
the Cartesian tree, so we build that tree once and answer each query as a lowest-common-ancestor
computation, pricing it at $\bigO{(n+q)\log n}$ -- acceptable here and structurally illuminating.

\solstep{Key Observations}
\begin{observation}
By \thmref{c26-rmqlca}, $\min a[a..b] = a[\lca(a,b)]$ in the min-Cartesian tree, so the entire problem
reduces to answering LCA queries on a tree we can build in $\bigO{n}$.
\end{observation}
\begin{observation}
LCA via binary lifting needs only the parent array (free from the build, \coderef{c26-stackbuild}) and
depths; the lifting table costs $\bigO{n\log n}$ and each LCA $\bigO{\log n}$.
\end{observation}
\begin{observation}
No updates occur, so the frozen tree is valid for every query -- the exact regime where the reduction
beats a stack pass that could not remember edits.
\end{observation}

\solstep{Think It Through With Me}
\begin{thinkbox}
\deadend{Scan $[a,b]$ for each query.}
{$\bigO{n}$ per query, $\bigO{nq}$ total -- $\sim4\cdot10^{10}$ at the limits, a \tle.}
\deadend{Segment tree.}
{Works ($\bigO{(n+q)\log n}$), but that is the engine of Chapter~18 (\chref{c14}); here the goal is to
\emph{use the Cartesian tree as an index}, so we take the LCA route instead -- it is the one that proves
\thmref{c26-rmqlca} does work.}
\aha{Read the query endpoints as tree nodes, lift to their LCA, output that node's value. The tree is
already built by the stack; lifting is the only extra machinery, and it is Chapter~16's, not ours.}
\end{thinkbox}

\solstep{Solution}
As in \coderef{c26-rmq}: monotonic-stack build to get \code{par, Lc, Rc} and \code{root}; iterative depth
sweep from \code{root}; binary-lifting table \code{up}; for each query compute $c=\lca(a,b)$ and print
$a[c]$. Because \thmref{c26-rmqlca} guarantees $a[c]$ is the range minimum, the answers are correct
regardless of the interior of the query.

\solstep{Proof of Correctness}
\begin{proofbox}[Each printed value is the range minimum]
For query $[a,b]$ let $m$ be a position of the minimum of $a[a..b]$. By \lemref{c26-cartmin}, $m=\lca(a,b)$.
The algorithm prints $a[\lca(a,b)]=a[m]$, which is by definition the minimum of the range. The LCA
computation itself is the textbook binary lifting, correct because lifting the deeper node to equal depth
and then ascending both in lock-step from the top of the table stops at the child of the LCA, whose
parent is returned. Hence every output equals the true range minimum. \end{proofbox}

\solstep{Implementation}
\begin{lstlisting}[caption={Static Range Minimum Queries -- Cartesian build plus binary-lifting LCA (CSES 1647)}, label={lst:c26-rmq1647-soln}]
#include <bits/stdc++.h>
using namespace std;
int main(){
    ios::sync_with_stdio(false); cin.tie(nullptr);
    int n, q; cin >> n >> q;
    vector<long long> a(n + 1);
    for (int i = 1; i <= n; ++i) cin >> a[i];

    vector<int> par(n+1,0), Lc(n+1,0), Rc(n+1,0), st;
    for (int i = 1; i <= n; ++i) {
        int last = 0;
        while (!st.empty() && a[st.back()] > a[i]) { last = st.back(); st.pop_back(); }
        if (!st.empty()) { Rc[st.back()] = i; par[i] = st.back(); }
        if (last)      { Lc[i] = last; par[last] = i; }
        st.push_back(i);
    }
    int root = st.front();

    vector<int> dep(n+1, 0);
    vector<int> fs = {root};
    while (!fs.empty()) {
        int u = fs.back(); fs.pop_back();
        if (Lc[u]) { dep[Lc[u]] = dep[u] + 1; fs.push_back(Lc[u]); }
        if (Rc[u]) { dep[Rc[u]] = dep[u] + 1; fs.push_back(Rc[u]); }
    }

    int LOG = 1; while ((1 << LOG) < n) ++LOG;
    vector<vector<int>> up(LOG, vector<int>(n+1, 0));
    for (int v = 1; v <= n; ++v) up[0][v] = par[v];
    for (int k = 1; k < LOG; ++k)
        for (int v = 1; v <= n; ++v) up[k][v] = up[k-1][ up[k-1][v] ];

    auto lca = [&](int u, int v) {
        if (dep[u] < dep[v]) swap(u, v);
        int d = dep[u] - dep[v];
        for (int k = 0; k < LOG; ++k) if ((d >> k) & 1) u = up[k][u];
        if (u == v) return u;
        for (int k = LOG - 1; k >= 0; --k)
            if (up[k][u] != up[k][v]) { u = up[k][u]; v = up[k][v]; }
        return par[u];
    };

    while (q--) { int A, B; cin >> A >> B; cout << a[ lca(A, B) ] << '\n'; }
    return 0;
}
\end{lstlisting}

\solstep{Trace}
\begin{tracebox}[Official CSES sample: array \texttt{3 2 4 5 1 1 5 3}]
The stack build (popping on strict $>$, so the two equal 1s are kept distinct) makes node $5$ the root,
with left subtree the block $1..4$ and right subtree $6..8$. Then: query $(2,4)$: node $2$ is an ancestor
of node $4$ (path $2\to3\to4$), so $\lca=2$, $a[2]=2$. Query $(5,6)$: $\lca(5,6)=5$ (root), $a[5]=1$.
Query $(1,8)$: $\lca(1,8)=5$ (root), $a[5]=1$. Query $(3,3)$: $\lca(3,3)=3$, $a[3]=4$. Output $2,1,1,4$
-- matches CSES. \checkmark
\end{tracebox}

\solstep{Complexity}
\begin{complexitybox}
Time $\bigO{(n+q)\log n}$: $\bigO{n}$ build, $\bigO{n\log n}$ lifting table, $\bigO{\log n}$ per query;
with $n,q\le2\cdot10^5$ this is a few million operations. Memory $\bigO{n\log n}$ for the table
($\approx 2\cdot10^5 \times 18$ integers, a few MB). Verdict: \ac. Note the sparse table (Chapter~19,
\chref{c13}) reaches $\bigO{1}$ per query with the same build -- use it when you need array RMQ alone; the
LCA reduction is the tool when the tree is already there.
\end{complexitybox}

\solstep{Pitfalls and Variants}
\begin{warningbox}
Watch the same three traps: query the \emph{positions} not the values; keep the depth sweep iterative so a
sorted input (a chain of depth $n$) does not overflow recursion; and do not re-derive LCA -- import it.
\emph{Variants:} for a range \emph{maximum} build the max-Cartesian tree (flip the pop to $<$) and the
identity becomes $\max = a[\lca]$. If updates appear, abandon the frozen tree for a segment tree
(\chref{c14}). If you need $\bigO{1}$ queries, replace the LCA table with a sparse table over $a$
directly (\chref{c13}) -- but then you have quietly left this chapter's reduction behind.
\end{warningbox}

\begin{insightbox}[Key takeaways]
\begin{enumerate}[nosep, leftmargin=1.6em]
  \item \textbf{A range minimum is a lowest common ancestor} -- $\min a[l..r]=a[\lca(l,r)]$
        (\thmref{c26-rmqlca}), a one-line corollary of the span lemma.
  \item \textbf{RMQ and LCA are the same problem} in two encodings; this section is
        array$\to$LCA, Chapter~16 is tree$\to$RMQ via the Euler tour.
  \item \textbf{The tree is an index, not just an image} -- the $\bigO{n}$ build \emph{plus} any LCA oracle
        is a working static-RMQ engine, which is how CSES 1647 falls (\probref{c26-rmq1647}).
\end{enumerate}
\end{insightbox}

\begin{hardeningbox}
\begin{itemize}[nosep, leftmargin=1.4em]
  \item \textbf{Second-smallest in a range} -- once you have the tree, the answer is the smaller of the
        two children along the $\lca$ path; a common contest twist on static RMQ.
  \item \textbf{Cartesian-tree DFS orders} -- inorder recovers the array; a preorder gives a "smallest
        first" enumeration used by some offline union tricks.
  \item \textbf{Range mode / nearest smaller} -- both read the same $L,R$ arrays from \secref{c26-monostack};
        the tree unifies them.
\end{itemize}
\end{hardeningbox}

\begin{recallbox}[10-second refresher -- 9.3 RMQ via the Cartesian Tree]
In the min-Cartesian tree the minimum of any contiguous range $[l,r]$ is stored at $\lca(l,r)$
(\thmref{c26-rmqlca}, from the span lemma \lemref{c26-cartmin}). So static RMQ reduces to LCA: build the
tree in $\bigO{n}$ (\coderef{c26-stackbuild}), record parents, lift to the LCA per query
(\coderef{c26-rmq}), print $a[\lca]$. The reverse reduction (LCA $\to$ RMQ on the Euler-tour depth array)
is owned by Chapter~16 (\chref{c11}); a pure array-only $\bigO{1}$-query RMQ is the sparse table of
Chapter~19 (\chref{c13}). CSES 1647 (\probref{c26-rmq1647}) is the worked reduction.
\selftest{
  \item Why is the range's minimum position an \emph{ancestor} of both endpoints?
  \item Why can no \emph{deeper} node be their common ancestor?
  \item State the RMQ $\equiv$ LCA equivalence and who owns which direction.
  \item Give the preprocessing and per-query costs of the lifting route.
  \item When should you prefer a sparse table to this reduction?
}
\end{recallbox}
